(1) 设反射光子能量为 \(E_r>0\)，镜子末态动量为 p、总能为 \(\mathcal E_m=\sqrt{M^2c^4+p^2c^2}\)。由于反射光子的动量为 \(-E_r/c\)，总能与动量守恒分别为
\[
Mc^2+E=\mathcal E_m+E_r,\qquad
E/c=p-E_r/c.
\]
所以 \(pc=E+E_r\) 且 \(\mathcal E_m=Mc^2+E-E_r\)。

(2) 将两式代入镜子的相对论质量壳关系：
\[
(Mc^2+E-E_r)^2-(E+E_r)^2=M^2c^4.
\]
化简得 \(2Mc^2(E-E_r)=4EE_r\)，因此
\[
E_r=\frac{E}{1+2E/(Mc^2)}.
\]
镜子动量 \(p=(E+E_r)/c\)，末速度可由 \(v=pc^2/\mathcal E_m\) 求出。令 \(\varepsilon=E/(Mc^2)\)，
\[
\frac vc=\frac{2\varepsilon(1+\varepsilon)}{1+2\varepsilon+2\varepsilon^2}<1.
\]
该结果是严格相对论反冲，不把镜子动能预先当作经典表达式。

(3) 镜子机械动能增加
\[
K_m=\mathcal E_m-Mc^2=E-E_r
=\frac{2E^2}{Mc^2+2E}\ge0.
\]
光子丢失的能量完全进入镜子动能，没有额外吸收或热激发；同时总动量也满足第一部分关系。

(4) 当 \(\varepsilon\ll1\)，反射能量
\(E_r/E=(1+2\varepsilon)^{-1}=1-2\varepsilon+O(\varepsilon^2)\)，故 \(\Delta E/E=(E_r-E)/E\simeq-2E/(Mc^2)\)，表现为反冲红移。若 M→∞，\(v\to0,K_m\to0,E_r\to E\)，恢复固定镜面对入射光子的理想等频反射。
